\chapter{Arquitectura e interfaces}

\section{Topología}

La topología activa separa presentación, API, datos y servicios externos:

\begin{lstlisting}
Navegador
  |-- HTTPS --> React/Vite en Vercel
  |               |-- /api/* --> Django/DRF/Daphne en Render
  |-- WSS --------------------> Django Channels en Render
                                  |-- PostgreSQL (Supabase)
                                  |-- Storage (Supabase)
                                  |-- Redis (si hay multiples procesos)
                                  `-- Brevo API o SMTP
\end{lstlisting}

El rewrite `/api` hace que la cookie de renovación sea first-party. WebSocket conecta directamente
a Render porque ese rewrite no actúa como proxy WebSocket.

\section{Responsabilidades por capa}

\begin{tabularx}{\textwidth}{@{}L{0.2\textwidth}Y Y@{}}
\toprule
\textbf{Capa} & \textbf{Responsabilidad} & \textbf{Regla de mantenimiento}\\
\midrule
React & páginas, componentes y estado de interfaz & componentes dependen de interfaces/contextos\\
Infraestructura FE & HTTP, WebSocket y configuración & una fuente tipada para variables y clientes\\
Views/serializers & HTTP, permisos y validación de entrada & no contienen consultas o reglas complejas\\
Services & casos de uso y reglas de negocio & una responsabilidad y dependencias explícitas\\
Repositories & consultas ORM & no exponen decisiones de presentación\\
Models/events & estado persistente e historial & cambios mediante migraciones versionadas\\
\bottomrule
\end{tabularx}

\section{API principal}

Las rutas no llevan barra final salvo donde se indica. La referencia exhaustiva está en los archivos
`urls.py` de cada aplicación.

\begin{longtable}{@{}L{0.16\textwidth}L{0.43\textwidth}L{0.31\textwidth}@{}}
\toprule
\textbf{Método} & \textbf{Ruta} & \textbf{Propósito}\\
\midrule
\endfirsthead
\toprule
\textbf{Método} & \textbf{Ruta} & \textbf{Propósito}\\
\midrule
\endhead
POST & /api/auth/register & registro de cliente\\
POST & /api/auth/login & acceso y cookie de renovación\\
POST & /api/auth/token/refresh & renovar sesión\\
POST & /api/auth/verify-email & confirmar correo\\
GET/POST & /api/tickets/ & listar por rol o crear como Cliente\\
GET & /api/tickets/\{id\} & detalle autorizado\\
GET & /api/tickets/\{id\}/historial & eventos del ticket\\
PATCH & /api/tickets/\{id\}/asignar & asignar como Administrador\\
PATCH & /api/tickets/\{id\}/reasignar & reasignar como Administrador\\
PATCH & /api/tickets/\{id\}/estado & cambiar estado como personal\\
POST & /api/tickets/\{id\}/comentario & añadir comentario\\
GET & /api/reportes/tickets & indicadores administrativos\\
POST & /api/reportes/exportar & exportar PDF o Excel\\
GET & /health/ & salud de aplicación y base\\
\bottomrule
\end{longtable}

\section{Autenticación y tiempo real}

El access token vive en memoria y se envía como Bearer. El refresh token viaja en cookie
`HttpOnly`, `Secure` en producción y `SameSite=Lax` con el rewrite first-party.

Los WebSocket disponibles son `/ws/notifications/` y `/ws/tickets/\{ticket\_id\}/`. El token de
acceso viaja como subprotocolo, no en la URL. El servidor valida origen, autenticación y acceso al
ticket antes de aceptar el canal.

\section{Datos y archivos}

{\hbadness=10000 PostgreSQL conserva usuarios, servicios, tickets, eventos, notificaciones y
reportes derivados. Supabase Storage guarda archivos cuando está configurado. La service role key
pertenece únicamente al backend. Los objetos, backups y exportaciones deben seguir la misma
política de retención y acceso que los registros relacionados.\par}

\RiskWarning{No exponga `SUPABASE\_SERVICE\_KEY` en variables `VITE\_*`, JavaScript, capturas,
logs o archivos de soporte.}
